% Standalone section chunks: include after the existing generating-function overview.
\section{计数建模：先确认哪些对象被区分}\label{knowledge-counting-model}
\index{计数建模}\index{OGF}\index{EGF}
写公式前，先用一句话说明：\textbf{一个方案是什么，什么变化会产生新方案？}
“把 $n$ 个物品分成 $k$ 组”不足以选公式；物品是否带标号、组是否带标号、
组内是否有序、组间是否有序、空组是否允许，分别决定不同问题。

\subsection{五个判别问题与两个最小反例}
\begin{enumerate}
\item \textbf{物品可区分吗？}给定标签 $1,\ldots,n$，交换两个物品可能改变方案，
属于带标号计数。相同球的方案只记录各类或各盒数量，不能再乘一个 $n!$。
\item \textbf{容器可区分吗？}盒子 A、B 是不同位置，即使容量相同也不同。
无标号盒子通常记录一个集合划分；空盒不形成新的块。
\item \textbf{次序在哪里？}一组元素组成集合、线性排列、圆排列，分别是不同结构。
例如固定 $s$ 个标签组成非空块时，集合只有 1 种，线性排列有 $s!$ 种，
不区分起点的圆排列有 $(s-1)!$ 种。不能只看到“无序分组”就忽略组内次序。
\item \textbf{哪些数量固定？}“恰好 $k$ 个非空组”与“至多 $k$ 个组”不同；
如果块数也要统计，可用额外变量 $y$ 标记，每个块贡献一个 $y$。
\item \textbf{空对象是否是一种方案？}通常空序列、空集合各有 1 种，所以常数项为 1。
“至少一个块”则要减去 1。无物品且无盒时通常计 1，非空物品放入零盒计 0；
题目另有定义时以题面为准。
\end{enumerate}

\textbf{不能机械除以盒子数的阶乘。}
两个不同球放入两个有标号且允许空的盒子，有 $2^2=4$ 种；
商去盒子标号后只有 $\{\{1,2\}\}$ 与 $\{\{1\},\{2\}\}$ 两种。
这个例子恰巧等于 $4/2!$，却不是通用理由：三个盒子时有 $3^2=9$ 种，
无标号且至多三块仍只有 2 种，而 $9/3!$ 甚至不是整数。
原因是含空盒时不同方案的稳定子大小可能不同。
若\emph{恰好 $k$ 个非空块}且只区分盒子标号，每个划分有恰好 $k!$ 个标号方式，才可除以 $k!$。

\textbf{不能把“顺序不同”当成“类型数量不同”。}
用长度 1、2 的砖铺长度 3 的一行，方案是 $111,12,21$，共 3 种；
若只问凑出总重量 3 的多重集合，则 $12$ 与 $21$ 是同一种，共 2 种。
前者是序列，后者是每种类型独立选择重数。

\subsection{OGF、EGF 的选择来自乘法，而非题目用词}
当独立部分只需把大小相加，且大小为 $i,j$ 的两个对象合并后恰有 $a_i b_j$ 种，
使用 OGF 的普通卷积。若还须从 $i+j$ 个不同标签中选 $i$ 个分给第一部分，
多出 $\binom{i+j}{i}$，使用 EGF 更自然。
\[
\text{OGF 数组： }A_i=a_i,\qquad
\text{EGF 数组： }A_i=a_i/i!,\qquad a_n=n!A_n.
\]
\textbf{卷积模板始终计算普通数组卷积；它不会识别并补上 EGF 的阶乘。}
题目出现标签不代表只能用 EGF，出现无序也不代表只能用 OGF；
判断依据是当前分解方式如何分配标签与次序。
一般非空带标号块的 EGF 记为 $B(x)$，则有序块序列是 $1/(1-B(x))$，
无序块集合是 $\exp B(x)$；二者都要求 $B(0)=0$，使每个固定次数只有有限项贡献。
本段组合解释参见 \href{https://next.oi-wiki.org/math/poly/egf/}{OI Wiki：指数生成函数}。

\section{组合数选型：先看模数，再看下标与预算}\label{knowledge-binomial-choice}
组合恒等式负责推导；本节负责把 $\binom nk$ 落到可运行的接口。
先用有符号或更宽类型处理负下标、$k>n$ 与 $n+k$ 溢出，再调用模板。
特别是 \texttt{Lucas::choose} 与 \texttt{ExLucas::choose} 接收无符号参数，
不要把容斥里的负数直接转换后传入。

\begin{description}
\item[素数 $p$，所有查询 $n\le N<p$：]
用第~\ref{compact-Binomial}~节（第~\pageref{compact-Binomial}~页）的 \texttt{Binomial<p> c(N)}，
查询 \texttt{c.choose(n,k)}；排列数用 \texttt{c.permute(n,k)}。
一次建表时间 $O(N+\log p)$、空间 $O(N)$，每次查询 $O(1)$。
后续用 \texttt{c.init(M)} 显式扩表；每次真正扩表都做一次扩展欧几里得求逆，
所以连续逐项扩表不能简单视作总共 $O(N)$。尽量一次扩至本批最大上界。
此处 $p$ 是编译期模数，调用者保证素数；模板的断言不等于判素。
\item[素数 $p$，$n$ 很大且 $p$ 足够小：]
用第~\ref{compact-Lucas}~节（第~\pageref{compact-Lucas}~页）的 \texttt{Lucas c(p)}，
再调用 \texttt{c.choose(n,k)}。
完整预处理时间 $O(p+\log p)$、空间 $O(p)$，单次 $O(1+\log_p n)$。
“下标大”不自动意味着适合此实现：$p=10^9+7$ 时，
两张长度为 $p$ 的 \texttt{int} 表约需 8 GB，不能靠 Lucas 定理消除这笔开销。
\item[合数 $m$，所有素数幂因子的表可以放下：]
用第~\ref{compact-ExLucas}~节（第~\pageref{compact-ExLucas}~页）的 \texttt{ExLucas c(m)}，
查询 \texttt{c.choose(n,k)}。它按 $m=\prod q_i$、$q_i=p_i^{e_i}$ 建表，
不是只预处理至最大查询的 $n$。空间 $O(\sum q_i)$，
建表时间可保守记为 $O(\sqrt m+\sum q_i+\sum\log q_i)$。
当前 \texttt{unit} 的每一层还有一次快速幂，故单次可保守记为
$O(\sum_i(1+\log_{p_i}(n+1))\log(n+2)+\sum_i\log q_i)$，
不能误写成纯粹 $O(\log n)$。模数 1 返回 0。
若某个 $q_i$ 很大，表就很大；ExLucas 不是任意 32 位合数的低内存黑盒。
\end{description}

\textbf{三者都不合适时不要硬套。}
若 $n$ 小、模数任意，Pascal 递推只用加法，算到第 $n$ 行耗时 $O(n^2)$、
滚动空间 $O(n)$；若只需 $k\le K$ 的列，可降至 $O(nK)$。
若素数 $p$ 很大、$k'=\min(k,n-k)<p$ 很小，
可将连续 $k'$ 个分子相乘，再乘 $(k'!)^{-1}$，耗时 $O(k'+\log p)$，
即使 $n\ge p$ 也可用，因为分母仍是单位；这需要自行写短循环。
合数模数下即便整数商存在，也不代表分母可逆，例如
$\binom42=6\equiv2\pmod4$，而 $2!$ 在模 4 下没有逆元。
不能把整数公式中的“除法”直接换成模逆元。

\section{生成函数建模到模板：三个可核验例子}\label{knowledge-generating-examples}
下面按“约束、因子、取系数、截断、算法”走完建模。
记 $N$ 为所求最大次数、$L=N+1$ 为数组长度，$M(L)=O(L\log L)$ 为 NTT 乘法复杂度。
第~\ref{compact-NttConvolution}~节（第~\pageref{compact-NttConvolution}~页）的 \texttt{NttConvolution<p,g>} 可指定素数与原根；
\texttt{FpsInverse}（第~\ref{compact-FpsInverse}~节，第~\pageref{compact-FpsInverse}~页）、
\texttt{FpsFunctions}（第~\ref{compact-FpsFunctions}~节，第~\pageref{compact-FpsFunctions}~页）、
\texttt{FpsPower}（第~\ref{compact-FpsPower}~节，第~\pageref{compact-FpsPower}~页）
当前固定为 $p=998244353$。
这些 FPS 操作的长度参数 $n$ 表示\emph{项数}，求 $x^N$ 必须传 $N+1$。

\subsection{例一：不同盒子、相同球、每盒有上下界（OGF 乘积）}
盒子 $i$ 要求 $l_i\le a_i\le u_i$，总数 $\sum a_i=N$。
令 $b_i=a_i-l_i$、$c_i=u_i-l_i$、$R=N-\sum l_i$。
若有 $u_i<l_i$ 或 $R<0$，答案为 0；否则所求为
\[
[x^R]\prod_{i=1}^k(1+x+\cdots+x^{c_i}).
\]
各因子选一个次数，唯一对应一个盒子的数量，所以系数不乘球的排列数。
例如容量为 $(2,3,1)$、下界均为 0、总数为 4，方案恰为
\[
(1,3,0),\ (2,2,0),\ (0,3,1),\ (1,2,1),\ (2,1,1),
\]
答案是 5。零容量的因子为 1，不应当删除这个盒子所携带的其他条件。

\textbf{截断。}只保留 $0,\ldots,R$ 次，因为因子没有负次数，高次不会回流到低次。
先把 $c_i$ 截至 $R$，每次乘法后仅当长度超过 $R+1$ 才截短；
不要把较短中间多项式补零至 $R+1$，否则会增加分治卷积的时间。
最终若数组尚无第 $R$ 项，该系数就是 0。
小规模用有界背包的滑动窗口：
\[
d_i(s)=d_i(s-1)+d_{i-1}(s)-[s>c_i]d_{i-1}(s-c_i-1),
\]
以 $d_0(0)=1$、负下标为 0 初始化，时间 $O(k(R+1))$、空间 $O(R+1)$，
只用加减，任意模数可用。不要为了出现乘积就立即写 NTT。

\textbf{卷积与特殊结构。}因子较多且总次数适中时，
分治调用 \texttt{NttConvolution<998244353>::multiply}；
设截短后的总次数为 $D=\sum\min(c_i,R)$，常用上界
$O(k+M(D+1)\log(k+1))$，每层截断通常还能减少工作。
若所有 $c_i=c$，令 $P(x)=1+\cdots+x^{\min(c,R)}$，
可用 \texttt{FpsPower::power(P,to\_string(k),R+1)}，时间 $O(M(R+1)+\log(k+1))$。
当 $k\ge1$，另一条路线是容斥：
\[
\sum_{j=0}^{\min(k,\lfloor R/(c+1)\rfloor)}
(-1)^j\binom kj\binom{R-j(c+1)+k-1}{k-1}.
\]
按组合数选型（第~\ref{knowledge-binomial-choice}~节，第~\pageref{knowledge-binomial-choice}~页）选择组合数接口，再比较项数与卷积成本。
$k=0$ 单独返回 $[N=0]$，不向公式塞入 $k-1=-1$。

\subsection{例二：步长 1、2 的有序序列（OGF 求逆）}
一次操作占用长度 1 或 2，序列总长度为 $N$，允许空序列。
一个操作的 OGF 为 $B(x)=x+x^2$；因为操作有先后，
\[
F(x)=\sum_{j\ge0}B(x)^j=\frac1{1-x-x^2},\qquad f_N=[x^N]F(x).
\]
$N=3$ 时有 $111,12,21$，$N=5$ 时为 8。
求前 $N+1$ 项可调用
\texttt{FpsInverse::inverse(\{1,-1,-1\},N+1)}
（初始化为其 \texttt{Poly} 类型），常数项为 1，可求逆。
若只处理这个二阶递推，$f_n=f_{n-1}+f_{n-2}$ 的 $O(N)$ 动态规划更快；
这里只是演示可推广到一般步长计数的建模，FPS 不必优于特化算法。

一般每个长度 $s\ge1$ 有 $b_s$ 种操作，则
$B(x)=\sum b_sx^s$，仍取 $(1-B)^{-1}\bmod x^{N+1}$，
时间 $O(M(N+1))$、空间 $O(N+1)$。
若 $N$ 巨大且允许步长有固定最大值，可证明有界阶递推后改用远项算法，
不应先分配长达 $N+1$ 的数组。

\textbf{反例。}若每种长度只记录使用多少次而不区分先后，应改为
$(1-x)^{-1}(1-x^2)^{-1}$，长度 3 的答案为 2。
若允许一种长度为 0 的操作，任意插入它便有无限多序列，
$B(0)=0$ 的条件被破坏；不能把一个本来不有限的计数问题交给求逆模板。

\subsection{例三：不同元素分成单点与无序对（EGF 指数）}
$n$ 个不同元素，每块大小为 1 或 2，块内与块间均无序。
固定的一个单点集合只有 1 种结构，一个二元集合也只有 1 种，故块的 EGF 为
\[
B(x)=x+\frac{x^2}{2!},\qquad
F(x)=\exp B(x),\qquad a_n=n![x^n]F(x).
\]
指数展开的 $B(x)^k/k!$ 先将标签分到 $k$ 个非空有序块，再去掉块的次序；
这里各块不空且标签集不交，除以 $k!$ 有严格的计数依据。
$n=4$ 时：全单点 1 种，一对加两单点有 $\binom42=6$ 种，
两对有 3 种，总计 10；不是系数 $[x^4]F=5/12$ 本身。

\textbf{落到当前接口。}在 $p=998244353$ 且 $N<p$ 下：
\begin{lstlisting}
using F = FpsFunctions;
using Z = F::Z;
F::Poly b(N + 1);
if (N >= 1) b[1] = 1;
if (N >= 2) b[2] = Z(2).inv();
auto f = F::exp(b, N + 1);
Binomial<998244353> comb(N);
Z answer = f[N] * comb.fac[N];
\end{lstlisting}
数组存的是除过阶乘的系数，最后还原；本例建模与 FPS 共耗时 $O(M(N+1))$。
若只求这一个模型，选出元素 1 的去向，还可独立得到
$a_0=a_1=1$、$a_n=a_{n-1}+(n-1)a_{n-2}$，给出 $O(N)$ 的简单算法与核验。

\textbf{增加块数条件。}恰好 $k$ 块改为
$n![x^n]B(x)^k/k!$。可调用
\texttt{FpsPower::power(b,to\_string(k),n+1)}，取第 $n$ 项后乘
\texttt{comb.fac[n]} 与 \texttt{comb.ifac[k]}，表预处理至 $\max(n,k)<p$。
先判 $k\le n\le2k$，不满足便为 0；空对象 $(n,k)=(0,0)$ 为 1。
例如 $(n,k)=(4,2)$ 只有三种两两配对。
若盒子本身带标号，就没有除以 $k!$ 的步骤；若盒子允许空，单盒因子还应添 1。

\subsection{抄代码前的四项检查与本地证据}
\begin{enumerate}
\item \textbf{目标系数与长度：}求 $[x^N]$ 使用长度 $N+1$；常数项是否代表空方案；
EGF 输入是否已乘逆阶乘，输出是否已乘阶乘。
\item \textbf{代数前提：}逆元分母非零；FPS 的 $\log$ 要求首项 1，$\exp$ 要求首项 0；
涉及阶乘和积分时不能跨过特征。整数有意义的计数，不保证其 EGF 在任意模数下可直接计算。
\item \textbf{实际容量：}NTT 接口检查的是\emph{完整乘积}长度 $a.size()+b.size()-1$，
不是最后保留的 $N+1$；“乘后截断”不会免掉中间数组。
固定模数的最大 NTT 长度为 $2^{23}$，上述 FPS 取 $L\le2^{22}$ 是安全的保守长度条件。
\texttt{FpsPower} 还显式检查这个半容量限制；不同 FPS 的局部乘积不能只凭最终次数估算。
\item \textbf{先用极小样例数对象：}至少检查空对象、一个元素、两个元素，
以及能区分“有序/无序”的第一个例子。错误模型也能被正确的卷积模板稳定算错。
\end{enumerate}

配套 \path{tests/math_knowledge_combinatorics.py} 用笛卡尔积枚举有界数量向量、
显式枚举步长序列、枚举集合划分，独立核对三个例子的系数与边界；
同时编译当前 C++ 接口并对照精确整数组合数。
运行命令为 \texttt{python3 tests/math\_knowledge\_combinatorics.py}。
这是本地建模与接口集成验证，不是新增在线评测 AC，也不取代各模板的完整回归。

\label{end-combinatorics-supplement}
